% C.3bt. Complete proof passed root and both peer audits on 19 November 2026.
\paragraph{C.3bt. Proper edge families (Proved).}
Let $k$ be algebraically closed of characteristic $p>0$.
Let $E=\bigcup E_i$ be a projective SNC exceptional surface in a
regular threefold, with connected double curves $C_e$ and individually
indexed triple points. Assume that $D=\sum a_iE_i$, $a_i>0$, satisfies
$\mathcal O_E(-D)$ ample. Put $S(E)=F_\gamma^2K_0(E)_{\mathbf Q}$,
and let $J_K(E)$ be the actual fixed smooth-proper-curve endpoint
relation group of EP. Then
\begin{equation}
 \ker\left(S(E)\longrightarrow
       \bigoplus_i\operatorname{CH}_0(E_i)_{\mathbf Q}\right)
                       \subseteq J_K(E).                 \tag{VE.1}
\end{equation}
There is no hypothesis on second dual cohomology, triple-point
Picard periods, or joint Picard restriction. In the exact resolved
threefold scope of C.3bm this removes the edge and triple-point kernel.
After C.3bo the remaining restrictions lie in the component Albanese
kernels, with their induced proper-family and exceptional relations.
If all component Albanese kernels vanish rationally, the ordinary
rational top quotient of the puncture is zero. The general component
kernel quotient is not computed here.

The new input is the following actual-family lemma. Its marked
restrictions are literally zero in $K_1$, not merely norm-zero.

\emph{Marked edge-family lemma.}
Let $C_1,\ldots,C_s$ be smooth connected proper curves over $k$,
with base points $c_e$ and finitely many marked points on each.
For $L_e\in\operatorname{Pic}^0(C_e)(k)$ and $u\in k^*$ there are
a smooth connected proper curve $P$, points $o,t\in P(k)$, and
actual classes $\xi_e\in K_1(C_e\times P)$ such that
\begin{align}
 \xi_e|_{C_e\times\{o\}}&=0,&
 \xi_e|_{C_e\times\{t\}}&=u([L_e]-1),\nonumber\\
 \xi_e|_{\{q\}\times P}&=0
 &&\text{for every marked point }q.                       \tag{VE.2}
\end{align}
The construction on smooth products is integral; the singular-scheme
application below uses rational coefficients.

\emph{1. A proper base and normalized divisors.}
If $u=1$ or every $L_e$ is trivial, take zero classes. Otherwise choose
a smooth connected projective curve
$A\subset\prod_e\operatorname{Pic}^0(C_e)$ through $0$ and the
nonzero tuple $(L_e)_e$. In dimension one use the whole abelian
variety; otherwise successive sufficiently high complete intersections
with prescribed points and first jets give such a curve in every
characteristic. The smoothness incidence estimate and connectedness
by duality and Serre vanishing are the same ones used in C.3bq
\citep[Lemma 33.47.3 and proof]{D_StacksBertiniUniformHigher}.
Denote its two points by $a_0,a_1$. The normalized Poincare bundles
$\mathcal M_e$ on $C_e\times A$ are trivialized along $c_e\times A$
and $C_e\times\{a_0\}$, and have fibres $L_e$ at $a_1$.

For every marked point $q$, put
$N_{e,q}=\mathcal M_e|_{\{q\}\times A}$, with its fixed
trivialization at $a_0$. Choose a rational section $r_{e,q}$ having
value $1$ in that trivialization at $a_0$, and divisor
$B_{e,q}=\operatorname{div}(r_{e,q})$ disjoint from $a_0,a_1$.
To do this, twist by a large effective divisor supported away from
the two points. Vanishing of $H^1$ after subtracting both points
makes evaluation onto, so prescribe value $1$ at $a_0$ and a
nonzero value at $a_1$, then divide by the twisting section.
Only finitely many divisors are involved.

\emph{2. Interpolate the covering function at their supports.}
Let $F$ be the union of all their supports. Choose a rational
function $h$ on $A$, with its only pole at $a_0$, satisfying
\begin{equation}
 h(a_1)=\frac{(u-1)^2}{u},\qquad h(x)=0\quad(x\in F),       \tag{VE.3}
\end{equation}
and having at least one simple zero away from $a_0,a_1$.
For the finite prescribed jet scheme $T$, choose $n$ large enough that
\[
                 H^1(A,\mathcal O_A(na_0)(-T))=0.
\]
This makes all the conditions compatible. A nonzero
leading polar coefficient at $a_0$ can be prescribed too.
One may impose simple zeros at every point of $F$ and at one
further point. This is the curve Riemann--Roch evaluation argument;
there is no bound on the characteristic. The value at $a_1$
is nonzero because $u\ne1$.

Normalize the double cover with function-field equation
\begin{equation}
             (g-1)^2=h g,\qquad
                    g^2-(h+2)g+1=0.                    \tag{VE.4}
\end{equation}
In odd characteristic its discriminant $h(h+4)$ has odd valuation
at a simple zero of $h$, proving irreducibility and separability.
In characteristic two, $y=(g+1)/h$ satisfies $y^2+y=1/h$.
The right side has a simple pole, whereas any pole of $z^2+z$
has even order, so this equation is again irreducible and separable.
Thus the normalization $f:P\to A$ is a finite separable degree-two
map of smooth connected proper curves.

Both $g$ and $g^{-1}$ are integral wherever $h$ is regular.
Every zero or pole of $g$ therefore lies above $a_0$.
Conversely $g+g^{-1}=h+2$ shows that $g$ is not a unit at any
point above that pole. Write $B=f^{-1}(a_0)_{\rm red}$.
Choose $o\in B$ and a point $t$ above $a_1$ with $g(t)=u$.
Such a point exists by (VE.4), including when it is ramified.
For every point $y$ above $F$, the same equation gives
\begin{equation}
                              g(y)=1.                    \tag{VE.5}
\end{equation}
This holds also in characteristic two and with ramification.

\emph{3. An actual Cech $K_2$ class across zeros and poles.}
Fix an edge and write $X=C_e\times P$,
$\mathcal M=(1\times f)^*\mathcal M_e$ and $V=C_e\times B$.
The bundle has its specified trivialization on $V$.
Choose a finite Zariski cover trivializing $\mathcal M$, using local
frames lifting that trivialization wherever an open meets $V$.
Such frames exist by locally lifting a basis of the invertible
module modulo the ideal of $V$ and shrinking.
Let $b_{ij}$ be the transition units, with the convention that
their Cech class is $c_1(\mathcal M)$.
Off $V$ take the symbols $\{g,b_{ij}\}$ in the Zariski sheaf
$\mathcal K_2$ of Quillen $K$-groups.

Near a point above $a_0$, choose a parameter $z$ on $P$ and write
$g=z^m v$, $b_{ij}=1+z b'_{ij}$, where $v,b_{ij}$ are units and
$m$ may be negative. In additive notation the extension is
\begin{equation}
        m\langle z,-b'_{ij}\rangle+\{v,1+z b'_{ij}\}.
                                                               \tag{VE.6}
\end{equation}
Indeed the Dennis--Stein convention is
$\langle r,s\rangle=\{r,1-rs\}$ when $r$ is invertible
\citep[Chapter III, Definition 5.11 and identities D1--D3,
pp.~43--44]{frWeibelKbookIII}.
Regular local Gersten injectivity identifies the sheaf $\mathcal K_2$
with its unramified subsheaf in the function-field group
\citep[Section 7, Theorem 5.11 and the preceding Gersten complex]
{rootQuillen1973}.
Thus the extensions glue uniquely and satisfy the Cech cocycle
identity, since they do so in the function field.
Changing compatible local frames gives a Cech coboundary: their
transition units are also $1$ modulo $z$, and the same formula
extends the corresponding symbols. We obtain
\begin{equation}
                         c_e\in H^1(X,\mathcal K_2).      \tag{VE.7}
\end{equation}
This uses regular local sheaf injectivity, not injectivity of an
arbitrary global affine $K_2$ group into its function field.
Functoriality of the displayed symbols gives the horizontal endpoint
and vertical marked-point restrictions used below.

At a zero or pole endpoint, (VE.6) specializes to
$m\langle0,-b'_{ij}|_{z=0}\rangle+\{v|_{z=0},1\}=0$;
the first zero follows directly from the defining Dennis--Stein
word. Hence $c_e|_o=0$. At $t$ it is the ordinary cup
$u\,c_1(L_e)\in H^1(C_e,\mathcal K_2)$.

\emph{4. Marked restrictions, including the boundary scalar.}
At a marked $q$, the restricted relative-trivialized bundle is
$f^*N_{e,q}$. The section $f^*r_{e,q}$ has value $1$ at every
point of $B$, and divisor
\[
 \sum_{x\in\operatorname{supp}B_{e,q}}\ \sum_{y\mid x}
                      e(y/x)n_x[y],
              \qquad B_{e,q}=\sum_x n_x[x].
\]
Use this section as frame away from its divisor. It is compatible
with the specified boundary trivialization. All remaining
transition calculations occur where $g$ is a unit.
The Gersten divisor formula for the restricted cup is consequently
\begin{equation}
 \sum_x\sum_{y\mid x} e(y/x)n_x\,g(y)[\mathcal O_y]
                          \quad\text{in }H^1(P,\mathcal K_2).
                                                               \tag{VE.8}
\end{equation}
A unit is here written additively as its $K_1(k)$ class.
This is also the proper point-pushforward projection formula
\citep[Section 7, Propositions 2.10 and 3.2]{rootQuillen1973}.
No boundary point-unit term is omitted: the comparison frame has
value $1$ at the boundary, so its comparison symbols extend by
(VE.6), with zero specializations. This explains the normalization
of the rational sections in Step~1.
Every $g(y)$ in (VE.8) equals $1$ by (VE.5), hence is zero in
$K_1(k)$. Ramification multiplicities, including wild ones in
characteristic two, make no difference. Every marked restriction
of $c_e$ is literally zero. No norm map is used.

\emph{5. Lift to the actual $K_1$ group by sheaf descent.}
For a noetherian finite-dimensional scheme there is the natural
strongly convergent descent spectral sequence
\begin{equation}
 E_2^{a,b}=H^a(X,\mathcal K_b)\Longrightarrow K_{b-a}(X),
 \qquad d_r:(a,b)\longmapsto(a+r,b+r-1).
                                                               \tag{VE.9}
\end{equation}
Use \citet[Theorem 10.3 and proof, pp.~383--386;
Appendix D.4--D.5 and proofs, pp.~424--426]{D_ThomasonTrobaugh1990}.
The filtration is the sheaf hypercohomology/Postnikov filtration,
not an assertion about pullback preserving codimension of supports.
Derived pullback of perfect complexes induces the map of $K$-presheaves
\citep[Sections 3.14--3.15, pp.~317--319]{D_ThomasonTrobaugh1990}.
The natural augmentation and functorial Postnikov construction give
the compatible descent-filtration maps for the horizontal and
vertical regular Cartier curves at issue.

As $\dim X=2$, no differential of order $r\ge2$ can enter or leave
$(a,b)=(1,2)$: its source would have negative cohomological degree
or its target would have degree at least three. Hence
\begin{equation}
 F^1_{\rm des}K_1(X)/F^2_{\rm des}K_1(X)
                              \simeq H^1(X,\mathcal K_2).
                                                               \tag{VE.10}
\end{equation}
Choose a lift $\xi_e\in F^1_{\rm des}K_1(X)$ of $c_e$.
For a regular curve the same descent filtration has $F^2_{\rm des}=0$
and $F^1_{\rm des}K_1\simeq H^1(\mathcal K_2)$.
Naturality makes the required restrictions of $\xi_e$ independent
of its choice and equal to the sheaf classes computed above.
The possible degree-two surface filtration ambiguity restricts
to zero on these Cartier curves, while choosing $F^1_{\rm des}$
excludes a determinant/unit term.
The ordinary cup on a smooth curve sends $u\,c_1(L_e)$ to
$u([L_e]-1)$. Thus (VE.2) holds for actual classes on the fixed
proper products. No formal or affine-only family has been substituted.

% Owned implication fragment for C.3bt.  The marked edge-family lemma is
% proved in the preceding cycles fragment.  No independent result row.
\subparagraph{From marked edge families to the whole component kernel.}
We use the notation and exact exceptional threefold scope of C.3bm and
C.3bo.  In particular $E=\bigcup_iE_i$ is a projective SNC surface over
an algebraically closed field of characteristic $p>0$, and a positive
exceptional combination $D=\sum_i a_iE_i$ satisfies
$\mathcal O_E(-D)$ ample.  Put
\[
 S=F_\gamma^2K_0(E)_{\mathbf Q},\qquad J=J_K(E),\qquad
 Q=\theta\operatorname{CH}_1(E)_{\mathbf Q},\qquad
 \varphi(s)=\bigl(\operatorname{ch}_2(s|_{E_i})\bigr)_i.
\]
Here $J$ is generated by the endpoint differences of actual perfect
gamma-two classes on $E$ times smooth connected projective curves, as
in C.3bm.  We prove, using the marked edge-family lemma above, that
\begin{equation}
 \varphi:S\twoheadrightarrow\bigoplus_i\operatorname{CH}_0(E_i)_{\mathbf Q},
 \qquad \ker\varphi\subset J.                         \tag{UK.1}
\end{equation}
Neither a dual-complex cohomology hypothesis nor a triple-point period
hypothesis is imposed.

\emph{1. Keep the actual totalization and prove pointwise surjectivity.}
For a fixed smooth projective curve $P$, or $P=\operatorname{Spec}k$,
let
\[
 \mathcal C_P^0=\bigoplus_iK(E_i\times P)_{\mathbf Q},\quad
 \mathcal C_P^1=\bigoplus_eK(C_e\times P)_{\mathbf Q},\quad
 \mathcal C_P^2=\bigoplus_vK(P)_{\mathbf Q},\quad
 \mathcal T_P=\operatorname{fib}(\mathcal C_P^1\to\mathcal C_P^2).
\]
The canonical geometric nullhomotopy, not a replacement chosen on
homotopy groups, gives
\begin{equation}
 K(E\times P)_{\mathbf Q}\simeq
          \operatorname{fib}(\mathcal C_P^0\to\mathcal T_P).
                                                               \tag{UK.2}
\end{equation}
This is the every-degree closed-union descent proved in (ET.6), from
\citet[Theorems A and C, printed pp.~29--30]{rootGeisserHesselholt2011}
and \citet[Theorem 8.1 and proof, p.~367]{D_ThomasonTrobaugh1990}.
The exact sequences, natural under both endpoint sections, are
\begin{align}
 \bigoplus_vK_2(P)_{\mathbf Q}&\longrightarrow\pi_1\mathcal T_P
 \longrightarrow\bigoplus_eK_1(C_e\times P)_{\mathbf Q}
 \longrightarrow\bigoplus_vK_1(P)_{\mathbf Q},             \tag{UK.3}\\
 \pi_1\mathcal T_P&\xrightarrow{\partial_P}K_0(E\times P)_{\mathbf Q}
 \longrightarrow\bigoplus_iK_0(E_i\times P)_{\mathbf Q}.
                                                               \tag{UK.4}
\end{align}

Every tuple of component zero-cycles has a unique pure weight-two class
\[
                         b\in\bigoplus_iK_0(E_i)_{\mathbf Q}.
\]
Its obstruction
$a\in\pi_0\mathcal T_k$ has eigenvalue $4$ for $\psi^2$ and maps to
zero in $\bigoplus_eK_0(C_e)_{\mathbf Q}$, whose weights are $0$ and $1$.
The first fibre sequence in (UK.2) therefore puts $a$ in the image of
$\bigoplus_vK_1(k)_{\mathbf Q}$, on which $\psi^2$ is multiplication by
$2$.  Consequently $4a=2a$, and $a=0$.  This uses only naturality of
the operations and exactness, not a semisimple weight decomposition of
$\pi_0\mathcal T_k$.  Equivalently that group is annihilated by
$(\psi^2-1)(\psi^2-2)^2$, which is invertible at the eigenvalue $4$.
Exactness lifts $b$ to $K_0(E)_{\mathbf Q}$; the rational weight-two
projector places the lift in $S$.  The finite gamma filtration and the
projectors on the projective surface and its products are supplied by
\citet[Chapter II, Proposition 4.9, Theorem 4.10, and Remark 8.8.5]
{D_Weibel2013KbookII}.  This proves the surjectivity in (UK.1).

\emph{2. Reduce every component-kernel class to degree-zero edge terms.}
Take $s\in\ker\varphi$.  The rational Chern character on the smooth
components makes its actual component $K_0$ restrictions zero, so
(UK.4) gives $y\in\pi_1\mathcal T_k$ with $\partial_k y=s$.
Replace $y$ by $(\psi^2-2)y/2$.  Its boundary remains $s$; on each
smooth proper double curve its edge image has weight two and is a
finite sum of terms $u([L]-1)$, with $u\in k^*$ and
$L\in\operatorname{Pic}(C_e)$.
This is the localization and Picard-times-units generation used in
C.3br and C.3bq
\citep[Section 7, Propositions 2.10 and 3.2]{rootQuillen1973}
\citep[Lemma 5.2 and its full proof, pp.~395--396]{D_PedriniWeibel1994}.
Commutation with the singular closed-union boundary uses the
arbitrary-scheme iterated-fibre and delooping construction of
\citet[p.~258, paragraph preceding Lemma 2.1]{frC3bcLevine1994};
we do not apply that lemma's regular-ambient hypothesis to $E$.

The weighted identity in (ET.9) makes alternating restriction onto
the edge degrees:
\[
 \bigoplus_i\operatorname{Pic}(E_i)_{\mathbf Q}
                  \twoheadrightarrow\bigoplus_e\mathbf Q.
\]
For each of the finitely many units, multiply component
divisor lifts of the edge degree vector by that constant unit in
$K_1$.  Subtract their canonical images in $\pi_1\mathcal T_k$.
Their boundaries in (UK.4) are zero.  The remaining edge image of $y$
is a finite sum of vectors
\begin{equation}
       \gamma_e=u([L_e]-1),\qquad
       L_e\in\operatorname{Pic}^0(C_e)(k),                \tag{UK.5}
\end{equation}
one unit $u$ at a time.  Clearing rational denominators before the
construction and dividing afterwards is harmless.  No surjectivity
of the full Picard restriction map has been assumed.

\emph{3. The marked families remove the remaining edge vectors.}
Apply the preceding marked edge-family lemma to (UK.5), with every
triple point on each connected edge included among the markings.
It supplies a smooth connected projective curve $P$ with endpoints
$o,t$ and actual classes
$\Gamma_e\in K_1(C_e\times P)_{\mathbf Q}$ such that
\begin{equation}
 (\Gamma_e)_t-(\Gamma_e)_o=\gamma_e,\qquad
 \Gamma_e|_{q_{e,v}\times P}=0\quad\hbox{in }K_1(P)_{\mathbf Q}
                  \quad\hbox{for every }v\supset e.       \tag{UK.6}
\end{equation}
The second equality is an equality on the actual proper parameter,
not just equality at its endpoints.  Thus (UK.3) gives
$\widetilde y\in\pi_1\mathcal T_P$ with edge image
$(\Gamma_e)_e$.  The class $\partial_P\widetilde y$ is an actual
perfect $K_0$ family.  On the three-dimensional projective scheme
$E\times P$, apply the rational Adams-weight-two projector to it.
Its endpoint difference therefore lies in $J$ and has the same
weight-two component-kernel contribution as
$\partial_k(\widetilde y_t-\widetilde y_o)$.

The difference between the originally chosen $y$ and these endpoint
lifts has zero edge image.  Exactness of (UK.3), with $P=k$, places
it in the image of $\bigoplus_vK_2(k)_{\mathbf Q}$.  The canonical
double boundary of this image is contained in $J$: the actual
proper-curve symbol construction in C.3bq Step~2 realizes every
Milnor symbol as a difference of evaluations of a class in $K_2$ of
a smooth proper curve, and naturality of (UK.3)--(UK.4) produces the
required family.  Its weight-two projection retains the desired
endpoint class.  This treats the full coherent-lift ambiguity;
no unique lift in (UK.3) is asserted or required.  Finite sums of
(UK.5) now prove $s\in J$, establishing (UK.1).

\emph{4. The precise component Albanese-kernel quotient.}
For each component write
\[
 T(E_i)=\ker\bigl(\operatorname{CH}_0(E_i)^0_{\mathbf Q}
       \longrightarrow\operatorname{Alb}(E_i)(k)\otimes\mathbf Q\bigr),
 \qquad T_E=\bigoplus_iT(E_i).
\]
Let $S_{00}$ have the meaning of C.3bo: both all component degrees
and all component Albanese images vanish.  Surjectivity in (UK.1)
gives $\varphi(S_{00})=T_E$.  Its kernel is contained in
$J\cap S_{00}$.  The exact denominator splitting of (AM.4) therefore
gives
\begin{equation}
 B^2(Z\setminus E)_{\mathbf Q}\simeq
 \frac{T_E}{
   \varphi(J\cap S_{00})+\varphi(Q\cap S_{00})}.          \tag{UK.7}
\end{equation}
The two denominator terms retain the induced actual proper-perfect
families and the actual exceptional cycle image, respectively.
For completeness, (AM.4) uses the Poincare subgroup
$U\subset J\cap Q\cap S_0$ whose component Albanese map is onto:
subtracting $u\in U$ first removes the Albanese coordinate, and for
$j+q\in S_{00}$ the replacement
$j+q=(j+u)+(q-u)$, with $\alpha_E(u)=\alpha_E(q)$, separates the
two denominator summands inside $S_{00}$.  Thus (UK.7) is not an
informal quotient by a sum of separately projected coordinates.

In particular, if all $T(E_i)$ vanish, the ordinary rational quotient
vanishes with no condition on $H^2(\Delta,\mathbf Q)$ or on triple
periods.  For arbitrary components the remaining group is exactly
(UK.7).  No assertion that all component Albanese kernels are induced
by globally compatible proper perfect families is made.


